\section{Trabalhos Relacionados}\label{sec:relacionados}

\cite{Graff2025} argumenta que, embora os métodos de XAI frequentemente sejam descritos como oferecendo ``razões" para as decisões de modelos de IA, essa caracterização é equivocada. O autor apresenta duas famílias de explicação por razão: causal/disposicional e interpretativa. Seu argumento é que nenhuma das duas famílias pode ser adequadamente aplicada a sistemas de IA. No primeiro caso, porque modelos de aprendizado de máquina não possuem atitudes normativas ou estados mentais que justifiquem suas decisões. No segundo, porque não é possível assumir que esses modelos possuam uma ``perspectiva" racional ou intencional que permita compreender suas ações em um ``sentido justificável". Assim, as explicações de XAI não são explicações por razões, mas sim explicações causais, que descrevem relações entre entradas e saídas sem envolver juízos normativos ou justificativos. O autor enfatiza as consequências práticas desse equívoco conceitual: ao apresentar as saídas de XAI como se fossem explicações de razões, corre-se o risco de criar falsas expectativas e uma confiança indevida nos sistemas de IA, levando usuários a acreditar que os modelos possuem justificativas racionais para suas recomendações. Em vez disso, defende que as explicações devem ser apresentadas em termos estritamente causais ou funcionais, evitando termos intencionais que sugerem agência ou racionalidade onde ela não existe. Esse diagnóstico reforça a necessidade de métodos que explicitem formalmente qual tipo de explicação está sendo gerado, em vez de tratarem todas as saídas de XAI como equivalentes.

Dando continuidade a essa discussão conceitual, o estudo de \cite{Pawlicki2024} aprofunda a lacuna epistemológica sob uma perspectiva metodológica, mostrando que a fragmentação não se limita ao conceito de explicação, mas também atinge os próprios critérios utilizados para avaliá-la. Através de uma revisão sistemática conduzida com rigor metodológico (PRISMA), os autores identificam um cenário marcado pela proliferação de propriedades, métricas e taxonomias que pretendem avaliar explicações em IA, mas que carecem de uniformidade terminológica, fundamentação conceitual comum e compatibilidade operacional. O estudo revela a existência de quase noventa métricas distintas utilizadas na literatura para mensurar ``qualidade explicativa", muitas delas redundantes, parcialmente sobrepostas ou introduzidas de forma idiossincrática por determinados grupos de pesquisa, o que torna inviável a comparação sistemática entre métodos de XAI. A contribuição mostra, assim, que o campo carece não apenas de uma definição clara do que conta como explicação, mas também de critérios sistemáticos e normativamente justificados para avaliar quando uma explicação é ``boa", ``adequada" ou ``suficiente".

Em uma proposta para minimizar sobreposições semânticas e agrupar diferentes terminologias que descreviam conceitos similares, \cite{Nauta2023} apresentou 12 propriedades, denominadas Co-12, para avaliar um sistema XAI. As Co-12 foram elencadas a partir de uma ampla revisão conceitual sobre qualidade de explicações, cobrindo tanto métodos \textit{ante-hoc} quanto \textit{post-hoc}, sendo o resultado um conjunto de propriedades categorizadas em três dimensões: conteúdo (\textit{correctness, completeness, consistency, continuity, contrastivity, covariate complexity}), apresentação (\textit{compactness, composition, confidence}) e usuário (\textit{context, coherence, controllability}). Somando a esse diagnóstico, \cite{Zhou2021} sistematiza métodos e métricas para avaliar explicações em ML e propõe uma base conceitual útil para comparar técnicas. O artigo distingue propriedades de interpretabilidade e de fidelidade, organiza explicadores em três classes (baseados em modelo, em atribuição e em exemplos) e descreve três eixos de avaliação. Um achado particularmente importante é o desbalanceamento no modo como a literatura avalia os explicadores: métodos baseados em modelo e em exemplos tendem a privilegiar critérios ligados à parcimônia e simplicidade estrutural, enquanto métodos baseados em atribuição concentram-se na fidelidade e na robustez. Essa assimetria ilustra como ainda há uma lacuna entre os esforços normativos, que descrevem o que deveria ser explicado, e os esforços técnicos, que medem o que de fato está sendo explicado, justificando a necessidade de uma visão integrada capaz de situar as propriedades de XAI dentro de estruturas explicativas consistentes em nível científico e não apenas operacional.

No contexto da interpretabilidade mecanicista em redes neurais com ativação ReLU, \cite{Barua2026} apresentam o método \textit{ReLU Region Reason} (Re3), que parte da propriedade de que essas redes se comportam de modo linear dentro de regiões específicas do espaço de entrada. Assim, para cada amostra, o método reescreve exatamente a computação da rede como um único mapeamento afim, \(f(x)=A_r x + D_r\), determinado pelos neurônios que permanecem ativos naquele caso. Isso permite calcular, de forma determinística e sem aproximações estocásticas, a contribuição exata de cada \textit{feature} nos \textit{logits} e nas probabilidades, além de identificar as regiões lineares em que a rede organiza o espaço de entrada. Os autores mostram alta concordância com LIME e SHAP nos rankings de importância, mas com estabilidade perfeita entre execuções. A relevância do Re3 está em mostrar que a geometria interna de redes ReLU pode ser explorada diretamente como mecanismo explicativo, aproximando-se de uma explicação causal-mecanicista ao reconstruir os caminhos computacionais efetivos de cada decisão.

\cite{Wickstrom2023} defende o método \textit{Representation Learning Explainability} (RELAX) como o primeiro método de explicação baseada em atribuição voltado especificamente para representações, e não para predições. A ideia central consiste em mascarar regiões da entrada e medir, via similaridade cosseno no espaço latente, o quanto a representação se altera, atribuindo importância às regiões cuja oclusão provoca maior mudança. O método é fundamentado teoricamente por uma interpretação em termos de \textit{kernels}: a importância de cada pixel corresponde a uma função de pontuação linear entre a representação original e a média ponderada das representações mascaradas no RKHS induzido pela função de similaridade. Além disso, o método incorpora quantificação de incerteza nas próprias explicações, permitindo filtrar regiões cuja atribuição é instável e aumentando a fidelidade do mapa de relevância resultante. A conexão explícita entre espaços de \textit{kernels} e atribuição de importância constitui um precedente direto para o presente trabalho, que estende essa lógica ao compor, por meio do produto tensorial de RKHS, matrizes de Gram construídas a partir de múltiplas fontes de um mesmo modelo, das quais os pesos, as ativações e os gradientes utilizados nos experimentos são apenas uma instanciação.

De forma complementar, \cite{Mae2026} propõem um arcabouço baseado em análise espectral para investigar a estrutura da qualidade de explicações em XAI. Os autores aplicam a decomposição em valores singulares (SVD) a uma matriz de redistribuição que codifica a atribuição de cada neurônio de saída a cada \textit{feature} de entrada, e demonstram que os valores singulares dessa matriz revelam dois fatores fundamentais da qualidade explicativa: estabilidade e sensibilidade ao alvo. A estabilidade é controlada pelo maior valor singular \(\sigma_1\) da matriz de redistribuição, que limita a amplificação de ruído durante o processo de explicação, enquanto a sensibilidade é capturada pela norma do espectro completo \(\|\sigma\|_2\), que mede a capacidade da explicação de distinguir entre diferentes conceitos na saída. Os autores mostram que esses dois fatores estão em tensão: parâmetros que aumentam a estabilidade tendem a reduzir a sensibilidade, e vice-versa, e propõem uma métrica combinada (SSM) que identifica o ponto de equilíbrio entre ambos. A relevância desse trabalho para o presente estudo reside em dois pontos. Primeiro, ele demonstra que a análise espectral é uma ferramenta poderosa para diagnosticar propriedades estruturais de explicações, o que corrobora a escolha da decomposição espectral de matrizes de Gram como operação central do arcabouço aqui proposto. Segundo, ao identificar estabilidade e sensibilidade como fatores espectralmente observáveis, ele oferece um vocabulário formal que se conecta à decomposição da energia espectral utilizada neste trabalho (Subseção \ref{subsec:spectral}), em que a energia retida nos modos dominantes, \(\widehat{E}_{\parallel}\), concentra as direções que organizam o espaço de representação, e a energia ortogonal, \(\widehat{E}_{\perp}\), isola aquilo que o subespaço retido não representa.

No campo da aprendizagem de representações, o trabalho de \cite{Dean2025} propõe um arcabouço matemático formal para descrever as transformações do mundo a partir da perspectiva de um agente em contextos de aprendizagem por reforço. Os autores partem do formalismo de representações \textit{disentangled} baseadas em simetria (SBDRL), que pressupõe que as simetrias do mundo devem estar presentes na representação interna do agente. Entretanto, demonstram que essa abordagem é restritiva, pois não consegue capturar transformações comuns em cenários de aprendizagem por reforço, como ações irreversíveis (por exemplo, consumir um item), ações restritas por obstáculos (paredes) ou mundos com ações não homogêneas. Para superar essas limitações, os autores generalizam tanto a condição de equivariância quanto a definição de \textit{disentanglement} do SBDRL utilizando teoria das categorias, estendendo-as de estruturas de grupo para qualquer álgebra, incluindo monoides e categorias pequenas. A relevância desse trabalho para o presente estudo reside na demonstração de que as propriedades de transformação do mundo podem ser formalizadas em estruturas algébricas mais gerais, o que se conecta diretamente com a proposta de utilizar a geometria dos espaços de representação como base para a interpretabilidade mecanicista, visto que tanto a abordagem algébrica quanto a geométrica buscam caracterizar as invariâncias e equivariâncias que organizam as representações internas dos modelos.

Os trabalhos revisados revelam três lacunas que o presente estudo pretende endereçar. Do ponto de vista conceitual, \cite{Graff2025} e \cite{Erasmus2021} mostram que o campo de XAI carece de clareza sobre qual tipo de explicação científica (dedutiva, indutiva, causal) está sendo gerado, e que essa distinção tem consequências práticas diretas para a confiança e a responsabilização. Do ponto de vista avaliativo, \cite{Pawlicki2024}, \cite{Nauta2023} e \cite{Zhou2021} evidenciam que as métricas existentes capturam dimensões distintas e por vezes incompatíveis da qualidade explicativa, sem um arcabouço unificador. Do ponto de vista técnico, os métodos geométricos e espectrais existentes, como o Re3 \citep{Barua2026}, o RELAX \citep{Wickstrom2023} e a análise espectral de redistribuição de \cite{Mae2026}, demonstram que a estrutura geométrica e espectral das representações internas contém informação explicativa rica, mas cada um deles opera sobre uma única fonte de informação e sobre uma construção específica, sem um arcabouço formal que garanta, para qualquer combinação de fontes, a existência de um espaço de representação com geometria e espectro bem definidos.

O presente trabalho responde a essas lacunas em dois níveis. No nível teórico, que constitui a sua principal contribuição, ele propõe o tensor de similaridade em RKHS e o pipeline núcleo-transformação, que compõem fontes heterogêneas de um mesmo modelo em um único espaço de representação e garantem, por demonstração, a semidefinição positiva, a pseudométrica e a decomposição da energia espectral sobre as quais as leituras são feitas, quaisquer que sejam as fontes e os \textit{kernels} escolhidos. Veja que essa construção é modelo-agnóstica, e os resultados mostram que ela se aplica, sem alteração, a um perceptron multicamadas, a uma rede convolucional, a um transformador e a uma árvore de decisão. No nível empírico, o trabalho responde à lacuna avaliativa com um protocolo em que cada leitura é formulada como uma pergunta sobre o comportamento do modelo (Subseção \ref{subsec:perguntas}), respondida contra uma referência externa à geometria e acompanhada de uma referência de nulidade construída para ela. Por fim, responde à lacuna conceitual declarando o tipo de explicação que cada leitura produz: as leituras de vizinhança, de erro e de confusão descrevem o comportamento do modelo, enquanto a leitura de diferenças entre uma amostra e o seu contrafactual e a datação são validadas por intervenção sobre a entrada ou sobre os estados internos, o que as aproxima da explicação causal-intervencionista e contrafactual discutida na Seção \ref{sec:referencial}.
